\documentclass{article}
\usepackage{amsmath, amssymb, amsthm}
\usepackage{hyperref}
\usepackage{geometry}
\geometry{margin=1in}

\title{Erd\H{o}s Problem \#222}
\date{Last updated: 2025-08-31}
\author{Source: erdosproblems.com}

\begin{document}
\maketitle

\noindent\textbf{Status:} OPEN \quad \textbf{No prize}

\noindent\textbf{Tags:} number theory, squares

\medskip

\noindent\textbf{Problem Statement:}

Let $n_1&#60;n_2&#60;\cdots$ be the sequence of integers which are the sum of two squares. Explore the behaviour of (i.e. find good upper and lower bounds for) the consecutive differences $n_{k+1}-n_k$.




Erd\H{o}s \cite{Er51} proved that, for infinitely many $k$,\[ n_{k+1}-n_k \gg \frac{\log n_k}{\sqrt{\log\log n_k}}.\]Richards \cite{Ri82} improved this to\[\limsup_{k\to \infty} \frac{n_{k+1}-n_k}{\log n_k} \geq 1/4.\]The constant $1/4$ here has been improved, most lately to $0.868\cdots$ by Dietmann, Elsholtz, Kalmynin, Konyagin, and Maynard \cite{DEKKM22}. The best known upper bound is due to Bambah and Chowla \cite{BaCh47}, who proved that\[n_{k+1}-n_k \ll n_k^{1/4}.\]The differences are listed at A256435 on the OEIS.




References

[BaCh47] Bambah, R. P. and Chowla, S., On numbers which can be expressed as a sum of two squares . Proc. Nat. Inst. Sci. India (1947), 101-103.

[DEKKM22] Dietmann, R. and Elsholtz, C. and Kaymynin, A. and Konyagin, S. and Maynard, J., Longer Gaps Between Values of Binary Quadratic Forms . International Mathematics Research Notices (2023), 10313–10349.

[Er51] Erd\H{o}s, P., Some problems and results in elementary number theory . Publ. Math. Debrecen (1951), 103-109.

[Ri82] Richards, Ian, On the gaps between numbers which are sums of two squares . Adv. in Math. (1982), 1-2.

\noindent\textbf{Related OEIS sequences:} A256435


\bigskip
\noindent\small{Source: \url{https://www.erdosproblems.com/222}}
\end{document}
